% Options for packages loaded elsewhere
\PassOptionsToPackage{unicode}{hyperref}
\PassOptionsToPackage{hyphens}{url}
\PassOptionsToPackage{dvipsnames,svgnames,x11names}{xcolor}
\PassOptionsToPackage{space}{xeCJK}
\documentclass[
  a4paper,
  10pt]{article}
\usepackage{xcolor}
\usepackage{amsmath,amssymb}
\setcounter{secnumdepth}{5}
\usepackage{iftex}
\ifPDFTeX
  \usepackage[T1]{fontenc}
  \usepackage[utf8]{inputenc}
  \usepackage{textcomp} % provide euro and other symbols
\else % if luatex or xetex
  \usepackage{unicode-math} % this also loads fontspec
  \defaultfontfeatures{Scale=MatchLowercase}
  \defaultfontfeatures[\rmfamily]{Ligatures=TeX,Scale=1}
\fi
\usepackage{lmodern}
\ifPDFTeX\else
  % xetex/luatex font selection
  \setmainfont[]{Times New Roman}
  \setmonofont[]{Menlo}
  \setmathfont[]{STIX Two Math}
  \ifXeTeX
    \usepackage{xeCJK}
    \setCJKmainfont[]{Songti SC}
  \fi
  \ifLuaTeX
    \usepackage[]{luatexja-fontspec}
    \setmainjfont[]{Songti SC}
  \fi
\fi
% Use upquote if available, for straight quotes in verbatim environments
\IfFileExists{upquote.sty}{\usepackage{upquote}}{}
\IfFileExists{microtype.sty}{% use microtype if available
  \usepackage[]{microtype}
  \UseMicrotypeSet[protrusion]{basicmath} % disable protrusion for tt fonts
}{}
\setlength{\emergencystretch}{3em} % prevent overfull lines
\providecommand{\tightlist}{%
  \setlength{\itemsep}{0pt}\setlength{\parskip}{0pt}}
\usepackage[a4paper,top=24mm,bottom=25mm,left=23mm,right=23mm,headheight=15pt]{geometry}
\usepackage{amsmath,amssymb,mathtools,bm}
\usepackage{booktabs,longtable,array}
\usepackage{graphicx}
\usepackage{float}
\usepackage{caption}
\usepackage{enumitem}
\usepackage{fancyhdr}
\usepackage{titlesec}
\usepackage{mdframed}
\usepackage{xcolor}
\usepackage{xurl}
\usepackage{hyperref}
\graphicspath{{../../}}

\definecolor{PaperBlue}{HTML}{000000}
\definecolor{PaperRule}{HTML}{B7C6D5}
\definecolor{PaperShade}{HTML}{F4F7FA}
\definecolor{PaperText}{HTML}{000000}

\hypersetup{
  colorlinks=true,
  linkcolor=PaperBlue,
  urlcolor=PaperBlue,
  citecolor=PaperBlue,
  pdfborder={0 0 0}
}

\color{PaperText}
\linespread{1.08}
\setlength{\parindent}{1.5em}
\setlength{\parskip}{0.25em}
\setlength{\emergencystretch}{3em}
\setlist{itemsep=0.2em,topsep=0.45em}
\setcounter{secnumdepth}{3}
\setcounter{tocdepth}{3}

\titleformat{\section}
  {\Large\bfseries\color{PaperBlue}}
  {\thesection}{0.7em}{}
\titleformat{\subsection}
  {\large\bfseries\color{PaperText}}
  {\thesubsection}{0.65em}{}
\titleformat{\subsubsection}
  {\normalsize\bfseries\color{PaperText}}
  {\thesubsubsection}{0.6em}{}
\titlespacing*{\section}{0pt}{2.2ex plus 0.5ex}{0.9ex}
\titlespacing*{\subsection}{0pt}{1.8ex plus 0.4ex}{0.7ex}
\titlespacing*{\subsubsection}{0pt}{1.5ex plus 0.3ex}{0.5ex}

\captionsetup{
  font=small,
  labelfont=bf,
  textfont=it,
  justification=centering,
  skip=6pt
}

\pagestyle{fancy}
\fancyhf{}
\fancyhead[L]{\small\nouppercase{\leftmark}}
\fancyhead[R]{\small Selene's Blog}
\fancyfoot[C]{\small\thepage}
\renewcommand{\headrulewidth}{0.4pt}
\renewcommand{\footrulewidth}{0pt}
\fancypagestyle{plain}{
  \fancyhf{}
  \fancyhead[R]{\small Selene's Blog}
  \fancyfoot[C]{\small\thepage}
  \renewcommand{\headrulewidth}{0.4pt}
}

\renewenvironment{quote}
  {\begin{mdframed}[
      linewidth=0.7pt,
      linecolor=PaperRule,
      backgroundcolor=PaperShade,
      leftline=true,
      rightline=false,
      topline=false,
      bottomline=false,
      innerleftmargin=10pt,
      innerrightmargin=8pt,
      innertopmargin=7pt,
      innerbottommargin=7pt,
      skipabove=8pt,
      skipbelow=10pt
    ]\itshape}
  {\end{mdframed}}

\newenvironment{theorembox}[1]
  {\par\medskip\noindent\textbf{Theorem (#1).}\quad}
  {\par\medskip}

\newenvironment{lemmabox}[1]
  {\par\medskip\noindent\textbf{Lemma #1.}\quad}
  {\par\medskip}

\newenvironment{proofbox}
  {\par\medskip\noindent\textit{Proof.}\quad}
  {\hfill$\square$\par\medskip}

\AtBeginDocument{\allowdisplaybreaks[2]}
\usepackage{bookmark}
\IfFileExists{xurl.sty}{\usepackage{xurl}}{} % add URL line breaks if available
\urlstyle{same}
% fallback for those not using the hyperref driver hyperxmp:
\makeatletter
\@ifundefined{xmpquote}{\newcommand{\xmpquote}[1]{#1}}{}
\makeatother
\hypersetup{
  pdftitle={CS224N \textbar{} Word2Vec Model},
  pdfauthor={Selene},
  colorlinks=true,
  linkcolor={black},
  filecolor={Maroon},
  citecolor={Blue},
  urlcolor={black},
  pdfcreator={LaTeX via pandoc}}

\title{CS224N \textbar{} Word2Vec Model}
\author{Selene}
\date{August 3, 2026}

\begin{document}
\maketitle

\section{Word Representation}\label{word-representation}

\subsection{Signifier and Signified}\label{signifier-and-signified}

A word is a signifier. It represents a signified object in the real or
imagined world. A signifier, or the corresponding signified object, can
refer to something very broad, such as uncle and drink, but it can also
refer to something much more specific, such as uncle Zuko and coffee.
Therefore, word meaning is extremely important, and even a small change
may have a large influence.

\subsection{Independent Words and
Vectors}\label{independent-words-and-vectors}

The simplest way to represent words is to treat them as mutually
independent and unrelated entities. Mutually independent components are
usually represented by one-hot vectors, that is, standard basis vectors.
For example,

\[
v_{\mathrm{tea}}=[0,0,1,0,\ldots,0]^T,
\qquad
v_{\mathrm{coffee}}=[0,0,0,0,\ldots,1,\ldots,0]^T.
\]

However, the problem with one-hot representations is that they cannot
encode any similarity relation. The dot product of any two different
vectors is \(0\). In reality, however, words in a sentence are related
to each other. Therefore, we introduce the next approach.

\subsection{Vectors from Manually Annotated Discrete
Attributes}\label{vectors-from-manually-annotated-discrete-attributes}

This approach manually constructs a feature vector whose entries
describe the attributes of a word. Clearly, this is not a good solution.

\section{Distributional Semantics and
Word2Vec}\label{distributional-semantics-and-word2vec}

\subsection{Distributional Hypothesis}\label{distributional-hypothesis}

The distributional hypothesis states that the meaning of a word can be
inferred from the distribution of contexts in which it appears. This is
one of the most influential ideas in NLP.

\subsection{Co-occurrence Matrix and Document
Context}\label{co-occurrence-matrix-and-document-context}

To implement the idea of the distributional hypothesis, we can construct
a co-occurrence matrix, denoted by \(X\). The size of this matrix is
\(\lvert V\rvert\times\lvert V\rvert\), where \(X_{ij}\) denotes the
number of times word \(i\) and word \(j\) appear in the same context.

In a sentence, we can mark multiple windows of different sizes. Short
windows tend to encode syntactic properties, while long windows tend to
encode semantic properties or even topic-level properties.

The problem with this approach is that high-dimensional vectors are
difficult to handle in modern neural systems, and high-frequency words
receive too much weight.

\subsection{Word2Vec Model and Objective
Function}\label{word2vec-model-and-objective-function}

Word2Vec represents each word in a fixed vocabulary as a low-dimensional
vector whose dimension is much smaller than the vocabulary size. It
learns the values of word vectors so that they can be used to predict
the distribution of context words through a simple function. The model
introduced here is \textbf{Skip-gram Word2Vec}.

\textbf{Skip-gram Word2Vec}

Let \(C,O\) be a pair of unknown random variables, where \(C \in V\)
denotes the center word, and \(O \in V\) denotes an outside word
appearing in the context of the center word. We use \(c,o\) to denote
their concrete values.

Suppose the vocabulary is \(V\). Each word has two \(d\)-dimensional
vector representations:

\begin{itemize}
\tightlist
\item
  \(u_w\): the vector of word \(w\) when it is used as a context word;
\item
  \(v_w\): the vector of word \(w\) when it is used as a center word.
\end{itemize}

We arrange all word vectors into matrices

\[
U,V\in\mathbb{R}^{\lvert V\rvert\times d}.
\]

Given center word \(c\), the probability that the context word is \(o\)
is defined as

\[
p_{U,V}(o\mid c)=
\frac{\exp(u_o^\top v_c)}
{\sum_{w\in V}\exp(u_w^\top v_c)}.
\]

Here, \(u_o^\top v_c\) represents the similarity between two word
vectors. The larger the dot product is, the more likely the model
believes that word \(o\) appears in the context of word \(c\).

Softmax converts the dot-product scores of all words into a probability
distribution:

\[
p_{U,V}(\cdot\mid c)\in\mathbb R^{\lvert V\rvert}.
\]

The goal of training Word2Vec is to make the predicted probability
distribution close to the true distribution of context words for word
\(c\) in the corpus, namely the row of the word co-occurrence matrix
corresponding to \(c\).

So far, we have only defined the model. We still need to determine the
parameters \(U,V\) through training. Word2Vec uses the cross-entropy
loss to make the predicted conditional distribution close to the true
context distribution:

\[
\min_{U,V}
\mathbb E_{o,c}
\left[
-\log p_{U,V}(o\mid c)
\right].
\]

Here, \((o,c)\) denotes a context word and center word pair sampled from
the corpus. For each sample, the model computes the negative log
probability of the true context word \(o\) given the center word \(c\):

\[
-\log p_{U,V}(o\mid c).
\]

If the model assigns a larger probability to the true context word, the
loss becomes smaller. Therefore, the training process continuously
adjusts \(U,V\) to increase the predicted probability of real word
pairs.

\subsection{Estimating Word2Vec from a
Corpus}\label{estimating-word2vec-from-a-corpus}

\textbf{Empirical Loss of Word2Vec}

Let \(D\) be a collection of documents \(\{d\}\). Each document is a
word sequence \(w_1^{(d)},\dots,w_m^{(d)}\), where \(w^{(d)} \in V\).
Let \(k \in \mathbb N_{++}\) be a positive integer window size. We now
connect the random variables \(C,O\) with this concrete dataset.

The center word \(C\) takes each word \(w_i\) in the document in turn.
For each such \(w_i\), the outside words are

\[
\{w_{i-k},\dots,w_{i-1},w_{i+1},\dots,w_{i+k}\}.
\]

Therefore, the objective function becomes

\[
L(U,V)=
\sum_{d\in D}
\sum_{i=1}^{m}
\sum_{\substack{j=-k\\j\ne0}}^{k}
-\log p_{U,V}
\left(
w_{i+j}^{(d)}
\mid
w_i^{(d)}
\right).
\]

Here, we sum over all documents, all words in each document, and all
words in the window. In this way, we accumulate the negative
log-likelihood of outside words given center words.

\textbf{Gradient-based Estimation}

We first make an initial guess with very little information for \(U,V\),
and then repeatedly move in the direction that locally improves this
guess the most. This gives a gradient-based implementation.

For a scalar function \(f\), its gradient with respect to the parameter
matrix \(U\) is denoted by \(\nabla_U f\). The gradient gives the
direction in which the function increases fastest. Therefore, if we want
to minimize the loss function, we should update the parameters in the
opposite direction of the gradient.

Before training starts, the word vector matrices \(U,V\) are usually
initialized as small random numbers:

\[
U^{(0)},V^{(0)}\sim N(0,0.001)^{\lvert V\rvert\times d}.
\]

In other words, each entry of the matrices is independently sampled from
a normal distribution with mean \(0\) and small variance. Then gradient
descent is used to update the parameters repeatedly. For example, the
update formula for matrix \(U\) is

\[
U^{(i+1)}=
U^{(i)}-
\alpha\nabla_U L\left(U^{(i)},V^{(i)}\right).
\]

Here, \(\alpha\) is the learning rate, which controls the step size of
each parameter update. Matrix \(V\) is updated in the same way:

\[
V^{(i+1)}=
V^{(i)}-
\alpha\nabla_V L\left(U^{(i)},V^{(i)}\right).
\]

By continuously updating \(U,V\) in the direction that decreases the
loss function, the predicted distribution of the model gradually
approaches the true distribution of context words.

\textbf{Stochastic Gradient}

For now, we do not consider how to compute the gradient explicitly.
Computing \(L(U,V)\) is very expensive because it requires traversing
the entire training set. Therefore, we use stochastic gradient
optimization.

\subsection{A Complete Gradient
Derivation}\label{a-complete-gradient-derivation}

First, we write down the gradient, and then use the linearity of
differentiation to move the gradient operator into the summation:

\[
\nabla_{v_c}\widehat L(U,V)=
\sum_{d\in D}
\sum_{i=1}^{m}
\sum_{\substack{j=-k\\j\ne 0}}^{k}
-\nabla_{v_c}
\log p_{U,V}
\left(
w_{i+j}^{(d)}
\mid
w_i^{(d)}
\right).
\]

Now we rewrite \(w_{i+j}^{(d)}\) as \(o\), and rewrite \(w_i^{(d)}\) as
\(c\):

\[
\begin{aligned}
\nabla_{v_c}\log p_{U,V}(o\mid c)
&=
\nabla_{v_c}\log
\frac{\exp(u_o^\top v_c)}
{\sum_{w\in V}\exp(u_w^\top v_c)} \\
&=
\nabla_{v_c}\log\exp(u_o^\top v_c)
-
\nabla_{v_c}\log
\sum_{w\in V}\exp(u_w^\top v_c).
\end{aligned}
\]

We take the gradient of the two parts with respect to \(v_c\)
separately.

The first part is

\[
\nabla_{v_c}\log\exp(u_o^\top v_c)=
\nabla_{v_c}(u_o^\top v_c)=
u_o.
\]

The second part is

\[
\begin{aligned}
\nabla_{v_c}
\log\sum_{w\in V}\exp(u_w^\top v_c)
&=
\frac{
\sum_{x\in V}
\exp(u_x^\top v_c)u_x
}{
\sum_{w\in V}
\exp(u_w^\top v_c)
} \\
&=
\sum_{x\in V}
p_{U,V}(x\mid c)u_x.
\end{aligned}
\]

Therefore,

\[
\nabla_{v_c}\log p_{U,V}(o\mid c)=
u_o-
\sum_{x\in V}
p_{U,V}(x\mid c)u_x.
\]

Here, \(u_o\) is the actually observed context word vector, while

\[
\sum_{x\in V}p_{U,V}(x\mid c)u_x
\]

is the expectation of the context word vector under the model's
predicted distribution. Therefore, this gradient can be understood as
observation minus expectation.

For the negative log-likelihood loss, the gradient has the opposite
direction:

\[
\nabla_{v_c}\left[-\log p_{U,V}(o\mid c)\right]=
\sum_{x\in V}
p_{U,V}(x\mid c)u_x-u_o.
\]

Gradient descent makes the center word vector \(v_c\) closer to the true
context word vector \(u_o\), while pushing it away from word vectors to
which the model incorrectly assigns high probabilities.

\end{document}
